\section{Билет 32. Внешняя политика России в начале XXI века}

\subsection{Внешняя политика России в 2000--2012 гг.}
К началу тысячелетия российская внешняя политика завершила стадию своего становления.

2000 г. "--- появление трёх новых важных документов в результате осмысления роли и места РФ в мировом сообществе:

\begin{enumerate}
\item Концепция национальной безопасности "--- содержит анализ внешних угроз интересам Российского государства.
\item Военная доктрина "--- развивает положения этой Концепции применительно к оборонному строительству.
\item Концепция внешней политики "--- решает ту же задачу применительно к конкретным областям внешнеполитической деятельности страны.
\end{enumerate}

апрель 2002 г. "--- в послании В.~В.~Путина Федеральному Собранию изложена суть внешней политики России:

\begin{itemize}
\item построение конструктивных нормальных отношений со всеми государствами мира;
\item необходимость для России быть сильной и конкурентоспособной в условиях жёсткой конкуренции за рынки, инвестиции, политическое и экономическое влияние;
\item прагматичное построение российской внешней политики, исходя из наших возможностей и интересов: военно"=стратегических, экономических, политических;
\end{itemize}

Реализуя идею разновекторной политики, Россия в 2000--2004 гг. поддерживала связи со многими странами мира на разных континентах.

\subsection{Теракт 2001 г. и ввод войск США и их союзников в Афганистан}

11 сентября 2001 г. "--- серия из четырёх террористических актов в США, в т.ч. столкновение угнанных террористами самолётов с башнями"=близнецами Всемирного торгового центра в Нью"=Йорке. Эти события привели к консолидации мирового сообщества в борьбе против мирового терроризма.

Октябрь 2001 г. "--- США начали крупную акцию возмездия в Афганистане, где, по их мнению, находился центр главной <<виновницы>> американской трагедии 11 сентября "--- руководимой бен Ладеном организации <<Аль"=Каида>>.

\subsection{Свержение режима Каддафи в Ливии}
Каддафи "--- лидер Ливии, отстаивавший справедливую политическую позицию и способствовавший процветанию ливийского народа. В феврале 2011 г. в Ливии начались массовые революционные выступления против режима Каддафи, которые переросли в гражданскую войну. В конце октября того же года Муаммар Каддафи был пойман и убит одной из вооружённых группировок в районе города Сирт. Запад не желал видеть Каддафи у власти, оказав поддержку ливийской оппозиции. Это привело к кровопролитной войне и массовому оттоку беженцев в Европу. Оказание поддержки <<арабской весне>> стало ошибкой с европейской стороны, ведь негативный эффект от событий в Ливии оказался намного более крупным, чем выгода, которую Запад извлёк для себя от уничтожения режима Каддафи.

\subsubsection{Позиция России по отношению к интервенции стран НАТО в Ливию}
Россия выступает в поддержку политического процесса комплексного урегулирования в этой стране, который должен осуществляться самими ливийцами, внешнее вмешательство в их дела считает контрпродуктивным. Также Россия считает необходимым последовательный вывод всех неливийских вооружённых группировок из страны.

\subsection{Позиция России по отношению к англо"=американскому вторжению в Ирак в 2003 г.}
20 марта 2003 г. "--- начало антитеррористической операции США и Великобритании в Ираке, формальным поводом для которой стала неподтверждённая информация о наличии у иракского руководства оружия массового уничтожения. 

Руководство России назвало агрессию <<серьёзной ошибкой США>>, эти события привели к росту антиамериканизма в мире. Помимо этого, усилилась дискуссия о месте ООН в структуре современных международных отношений, стали более интенсивными контакты России со странами СНГ, в ходе которых наибольшее внимание уделялось решению проблем безопасности, совместным экономическим проектам.

После вторжения прокатилась новая волна международного терроризма и экстремизма, это коснулось и России: 1 сентября 2004 г. в Беслане была захвачена школа, в результате действий боевиков погибли 333 человека.

\subsection{Гражданская война в Сирии}
Весной 2011 г. началось гражданское противостояние в Сирии, которое постепенно переросло в восстание против режима Башара Асада. В этот конфликт оказались вовлечены соседние страны, международные организации и мировые державы, в т.ч. и Россия. Ведущую роль в борьбе с правительством стали играть исламистские группировки: сирийское отделение <<Аль"=Каиды>> и <<Исламское государство>>.

\subsubsection{Позиция России по отношению к вмешательству США и их союзников в гражданскую войну в Сирии}
Кризис в Сирии приводит к распространению терроризма и экстремизма по всему миру, поэтому Россия предпринимает меры по нормализации ситуации в этой стране. В то же время Россия не приемлет применения принудительных мер и вооружённой силы в обход ООН, как это делает США, поскольку это способствует затягиванию противоречий на Ближнем Востоке. Действующую власть в Сирии Россия поддерживает как единственную легитимную.

\subsection{Продолжение расширения НАТО на восток. Вступление РФ в ВТО, ШОС и БРИКС}

Россия неоднократно возражала против расширения НАТО в направлении своей территории. Однако США продолжали активные контакты с Азербайджаном, Грузией и Украиной, которые выражали намерение вступить в альянс. Помимо этого США озвучили планы по размещению элементов системы противоракетной обороны на территории Чехии и Польши. Военное руководство России назвало систему ПРО <<серьёзным дестабилизирующим фактором, который может оказать значительное влияние на региональную и глобальную безопасность.

\subsubsection{ВТО}
Более 10 лет велись переговоры о присоединении России к Всемирной торговой организации (ВТО). 22 августа 2012 г. Россия стала полноправным членом ВТО с целью получения лучших условий для экспорта отечественных товаров, привлечения инвестиций извне и участия в формировании международных правил торговли с учётом национальных интересов. Вступление в ВТО открыло непростой период адаптации к новым внешним условиям конкуренции. В декабре 2005 г. президент выдвинул задачу превратить страну в лидера мировой энергетики, <<великую энергетическую державу>>.

\subsubsection{ШОС}
В 2006--2009 гг. стабильно во всех областях развивались отношения между Россией и Китаем. Здесь был оформлен региональный союз "--- Шанхайская Организация Сотрудничества (ШОС), куда помимо России и Китая вошли государства Центральной Азии, а на правах ассоциированных членов "--- Индия и Иран. Организация исключает участие США: этой стране было отказано даже в ассоциированном членстве.  

\subsubsection{БРИКС}
БРИКС (Brazil, Russia, India, China, South Africa) "--- межгосударственное объединение, которое было основано в июне 2009 г. в рамках Петербургского экономического форума. Его целью является развитие и укрепление экономического и инвестиционного взаимодействия.

\subsection{Создание ОДКБ}
15 мая 1992 г. "--- подписание договора о коллективной безопасности (ДКБ), который был зарегистрирован 1 ноября 1995 г. в Секретариате ООН. В 2009 г. в ОДКБ (<<Организация договора коллективной безопасности>>) входили Армения, Белоруссия, Казахстан, Киргизия, Россия и Таджикистан. В структуре ОДКБ основное внимание уделялось созданию мобильных коллективных сил быстрого реагирования, что связано с одной из главных задач организации "--- противодействия международному терроризму.
\subsection{Образование Союзного государства России и Белоруссии}

Республика Беларусь и РФ являются государствами"=участниками Договора о создании Союзного государства от 8 декабря 1999 г. Одним из важных ориентиров является создание единого экономического пространства. Сотрудничество России и Белоруссии эффективно развивается в таких сферах, как энергетика, транспорт, промышленная кооперация. В рамках Союзного государства страны также добились реализации единого таможенного пространства и зоны свободной торговли, граждане стран получили возможность пересекать границу по внутренним документам, была улучшена координация внешней политики.

\subsection{<<Цветные революции>>}
В результате внутренних трудностей и не без воздействия извне в ряде соседних стран сменилось руководство. Новые власти не только переориентировали свою политику на Запад, но и перешли на откровенно недружественные в отношении России позиции. Такие революции получили название <<цветных>> и были активно поддержаны США и странами Евросоюза.

2003 г. "--- <<революция роз>> в Грузии.

2004 г. "--- <<оранжевая революция>> на Украине.

2006 г. "--- <<революция тюльпанов>> в Киргизии.

Неудачей закончились попытки повторить этот сценарий в Узбекистане (2005 г.) и Белоруссии (2006 г.).

Всё это предъявило новые требования к внутренней устойчивости Российского государства и его внешней безопасности.

\subsubsection{Россия и <<оранжевая революция>> 2004 г. в Украине}

21 ноября 2004 г. ЦИК Украины объявила о победе на президентских выборах В.~Януковича. С победой его поздравили председатель Госдумы Б.~Грызлов и президент В.~Путин. В то же время большинство иностранных наблюдателей считали, что перевес Януковича в 3\% по сравнению с В.~Ющенко был достигнут за счёт нечестных выборов. В результате этого начался митинг на Майдане Незалежности в центре Киева, который проходил там непрерывно в течение двух месяцев: протестующие организовали там палаточный лагерь, а в некоторые дни митинги собирали до 1,5 млн. человек. Протестующие смогли добиться повторного голосования (эта идея была поддержана и европейскими лидерами как способ выхода из кризиса), в ходе которого победу одержал Ющенко. Оранжевая революция стала отправной точкой резкого крена украинской власти в сторону Запада.

\subsection{События 2008 г. в Южной Осетии}
8 августа 2008 г. "--- в Южной Осетии началась военная акция против <<сепаратистов>>, организованная в Грузии с ведома США. Война продолжалась 5 дней и завершилась разгромом грузинских войск. Это было результатом поддержки Южной Осетии со стороны России. Следствием Пятидневной войны стало признание Россией независимости Южной Осетии и Абхазии, с которыми были заключены договоры о дружбе и взаимопомощи. Такая позиция России привела к охлаждению отношений с Евросоюзом, который тогда безоговорочно поддержал Грузию.

\subsection{Принятие новой военной доктрины (2010)}
12 июля 2008 г. "--- Президент утвердил новую Концепцию внешней политики.

В 2010 г. Президентом РФ Д.~Медведевым была принята новая военная доктрина. Её наиболее важными положениями являются следующие:

\begin{itemize}
\item военная политика РФ направлена на недопущение гонки вооружений, сдерживание и предотвращение военных конфликтов;
\item ядерное оружие остаётся фактором сдерживания, но в то же время может быть применено как ответная мера по решению президента;
\item агрессия против одного из участников Союзного государства или ОДКБ рассматривается как агрессия против всех членов организации.
\end{itemize}
